% Plantoir house style: worked examples (the skill's section 4). A pattern to
% copy FROM, never a starter: copy a box into a handout made from
% starters/Handout.tex, whose preamble this file repeats line for line.
%
%   A. Substitute and evaluate: the substituted values in red (\sv).
%   B. A method step by step: steps aligned on =, a blue side note beside
%      the step that needs one (\sidenote in an alignat), a \method heading.
%   C. A table of values and its graph side by side INSIDE the box: two [t]
%      minipages, the right one starting with \vspace{0pt}; the curve has
%      no context, so it runs to the grid's edge with arrowheads; its label
%      sits on a peach patch.
%   D. Two short workings side by side inside one box (both intercepts).
%   E. A grey box of instructions before a group of questions (groupnote),
%      and how the answer key hides every example behind a \strand heading.
%   F. A small framed note beside an example (notebox).
%   G. A slope drawn on the graph: the run and the rise, dashed, from the
%      y-intercept (slopestep), for an example whose point is the slope.
%
% What makes an example good, in this house style: it teaches the WHOLE
% method a question below it needs, in the order a student does it; the
% value put in is marked; the step a student gets wrong carries a note
% beside it; a method that is graphical shows the table and the graph in
% the box, not only words. One example comes before each group of
% questions it serves. In the answer key the examples are hidden and the
% final answers given.
%
% ---------- The starter's preamble, line for line (a test keeps it so) ----------
\documentclass[11pt]{article}
\usepackage[letterpaper, left=0.75in, right=0.75in, top=1.05in, bottom=0.9in,
            headheight=14pt, headsep=0.3in]{geometry}
\usepackage{lmodern}
\usepackage[T1]{fontenc}
\usepackage{amsmath, amssymb}
\usepackage{xcolor}
\usepackage{enumitem}
\usepackage{tasks}
\usepackage{fancyhdr}
\usepackage{lastpage}
\usepackage{pgfplots}
\usepackage{array}
\usepackage{collcell}     % the C column's minus-sign warning (below)
\usepackage{needspace}
\usepackage[most]{tcolorbox}
\usepackage{nicematrix}   % tables: every rule drawn after the heading's shade (the skill's section 9)
\pgfplotsset{compat=1.18}

% ---------- Fill in ----------
% The school's name comes ONLY from print_school_name in the course's
% settings, or from the teacher in this conversation. Never a placeholder:
% with no name, leave the braces empty and the corner prints nothing.
\newcommand{\schoolname}{}
\newcommand{\coursecode}{COURSE}        % the course's code, in capitals
\newcommand{\handouttitle}{Handout Title}

% ---------- Questions, answers or hints ----------
\newif\ifanswers
\ifdefined\showanswers \answerstrue \fi
\newcommand{\qa}[2]{\ifanswers #2\else #1\fi}
\newcommand{\onlyq}[1]{\ifanswers\else #1\fi}
\newcommand{\onlya}[1]{\ifanswers #1\fi}
% The hints sheet (the skill's section 13): the hints alone, one per
% question, in the block right after \begin{document}. Nothing else in the
% file changes for hints, so the questions and the answers compile exactly
% as they would without them.
\newif\ifhints
\ifdefined\showhints \hintstrue \fi
\newcommand{\hint}[2]{\item[Question #1] #2}

% A blank students write on, in a question (the skill's section 4): an
% invisible strut makes its line about two lines tall, so a wrapped prompt
% still leaves a full line of writing room above the blank. \fillin{answer}
% or \fillin[5cm]{answer}: the student copy shows the blank, the answer key
% writes the answer on the same line, at the same height.
\newcommand{\fillin}[2][3.2cm]{\ \qa{\rule[-0.2em]{#1}{0.4pt}}{\makebox[#1]{#2}\llap{\rule[-0.2em]{#1}{0.4pt}}}\rule[-0.2em]{0pt}{2.1em}}

\ifanswers
  \newcommand{\doctitle}{\handouttitle\ --- Answers}
\else\ifhints
  \newcommand{\doctitle}{\handouttitle\ --- Hints}
\else
  \newcommand{\doctitle}{\handouttitle}
\fi\fi

% ---------- Header and footer (matches pages printed from the class website) ----------
\pagestyle{fancy}
\fancyhf{}
\renewcommand{\headrulewidth}{0pt}
\newcommand{\fillblank}[2]{#1 \rule[-0.15em]{#2}{0.4pt}}
% Lay the corners out as the course's print settings do (the skill's
% sections 2 and 5). The blanks are print_blanks' list in the order Name,
% Date, Class #; with no list, Name and Date, as here. Add
% \quad\fillblank{Class \#}{1.4cm} when the list has class_number.
\fancyhead[L]{\small\sffamily\color{gray}\schoolname}
\fancyhead[R]{\small\sffamily\color{gray}\fillblank{Name}{3.2cm}\quad\fillblank{Date}{2.2cm}}
\fancyfoot[L]{\small\sffamily\color{gray}\doctitle\ $\cdot$ \coursecode}
\fancyfoot[R]{\small\sffamily\color{gray}Page \thepage\ of \pageref{LastPage}}
\setlength{\parindent}{0pt}

% ---------- Colours and boxes ----------
\definecolor{why}{gray}{0.4}
\definecolor{peach}{RGB}{253,244,233}
\definecolor{peachframe}{RGB}{226,160,96}
\definecolor{exampletitle}{RGB}{176,38,30}
\definecolor{sub}{RGB}{200,30,30}      % a substituted value, in a worked example
\definecolor{note}{RGB}{20,70,140}     % a side note beside a step
\definecolor{tablehead}{RGB}{225,232,245}   % a table's heading row

% Worked example box
\newtcolorbox{example}[1]{
  enhanced, breakable=false,
  colback=peach, colframe=peachframe, boxrule=0.8pt, arc=3mm,
  left=3mm, right=3mm, top=2mm, bottom=2mm,
  before skip=0.9em, after skip=0.6em,
  fontupper=\normalsize,
  title={#1}, coltitle=exampletitle, colbacktitle=peach,
  fonttitle=\large\bfseries\sffamily,
  titlerule=0pt, toptitle=1mm, bottomtitle=-1mm,
}
% Inside a worked example: a step's method as a small heading, a note beside
% a step (in an alignat, after &&), and a substituted value in red.
\newcommand{\method}[1]{\medskip{\bfseries\sffamily #1}\par\smallskip}
\newcommand{\sidenote}[1]{\text{\small\sffamily\color{note}#1}}
\newcommand{\sv}[1]{{\color{sub}#1}}
% A small framed note set beside an example, in a [t] minipage
\newtcolorbox{notebox}{colback=white, colframe=peachframe, boxrule=0.5pt, arc=1.5mm,
  left=2mm, right=2mm, top=1mm, bottom=1mm, fontupper=\small\sffamily}
% A grey box of instructions before a group of questions
\newtcolorbox{groupnote}{colback=gray!8, colframe=gray!45, boxrule=0.6pt, arc=2mm,
  left=3mm, right=3mm, top=1.5mm, bottom=1.5mm, before skip=0.4em, after skip=0.8em,
  fontupper=\small}

% Section heading for the answer key, where the worked examples are hidden
\newcommand{\strand}[1]{\bigskip{\large\bfseries\sffamily #1}\par\smallskip}

% ---------- Tables (the skill's section 9) ----------
% Every value sits in a row \valuepitch tall, centred in it, so the first
% value is half a row below the heading rule and the last half a row above
% the bottom rule. The value columns of a table are one width
% (\setvaluecol{the widest value in either column} sets \valuecol to it plus
% 0.5cm). C is a centred value column, for a column with no negative value
% and not mixing whole numbers with decimals. N{$18$} is for every other
% column of numbers: each value right-aligned in a box as wide as the
% column's widest value WITHOUT its sign, the box centred, so the digits
% share one trailing edge and a minus sign juts out to the left.
% A C cell holding a minus sign prints a small red "use N" beside it, so a
% centred column of negatives shows on the rendered page (the skill's
% section 8). An N column never does. \renewcommand{\valuepitch}{3em} for a
% table with room to write.
\renewcommand{\arraystretch}{1.1}
\newcommand{\valuepitch}{1.5em}
\newcommand{\rowstrut}[1]{\rule[\dimexpr0.32em-0.5\dimexpr#1\relax\relax]{0pt}{#1}}
\newlength{\valuecol}
\newcommand{\setvaluecol}[1]{\settowidth{\valuecol}{#1}\addtolength{\valuecol}{0.5cm}}
\makeatletter
\newcommand{\centredvalue}[1]{%
  \begingroup\edef\x{\endgroup\noexpand\in@{\detokenize{$-}}{\detokenize{#1}}}\x
  \ifin@\else\begingroup\edef\x{\endgroup\noexpand\in@{\detokenize{$ -}}{\detokenize{#1}}}\x\fi
  \ifin@\llap{\tiny\sffamily\color{red}use N\ }\fi#1}
\makeatother
\newcolumntype{C}{>{\rowstrut{\valuepitch}\collectcell\centredvalue}w{c}{\valuecol}<{\endcollectcell}}
\newcolumntype{N}[1]{>{\rowstrut{\valuepitch}\settowidth{\dimen0}{#1}\hbox to\dimen0\bgroup\hss}w{c}{\valuecol}<{\egroup}}
% One heading-row height for tables set side by side
\newcommand{\headstrut}{\rule[-1.1em]{0pt}{2.6em}}

% ---------- Numbered questions ----------
\newenvironment{questions}
  {\begin{list}{}{\setlength{\leftmargin}{2.2em}\setlength{\labelwidth}{1.8em}
                  \setlength{\labelsep}{0.4em}\setlength{\itemsep}{0.7em}}}
  {\end{list}}
% Optional argument: lines to keep together so a question is never split
% from its parts. \roomfor is needspace's test with a neutral break: the
% package's own \needspace leaves a \penalty-100, and with a ragged bottom
% TeX then prefers to end the page THERE, before a question that would have
% fitted (measured on this starter: a page a quarter empty).
\makeatletter
\newcommand{\roomfor}[1]{\par\penalty\z@\begingroup\setlength{\dimen@}{#1}%
  \dimen@ii\pagegoal\advance\dimen@ii-\pagetotal
  \ifdim\dimen@>\dimen@ii\ifdim\dimen@ii>\z@\vfil\fi\penalty-\@M\fi\endgroup}
\makeatother
\newcommand{\q}[2][5]{\roomfor{#1\baselineskip}\item[\textbf{#2.}]}

\settasks{label=\alph*), label-width=1.4em, item-indent=1.8em, after-item-skip=0.35em}

% ---------- Graph styles (the skill's section 12) ----------
% Tick labels are drawn above the plotted lines on a white patch, so a line
% that passes a label never runs through its digits. The axes run a little
% past the grid, so a line's arrowhead at the grid's edge is clear of the
% axis arrows. Leave a tick label out where an arrowhead ends on it
% (yticklabels={,$-4$,...}).
\pgfplotsset{
  layers/labels on top/.define layer set={axis background, axis grid, axis ticks, axis lines,
    main, axis tick labels, axis descriptions, axis foreground}{},
  gridgraph/.style={
    set layers=labels on top,
    axis lines=middle, grid=both,
    major grid style={gray!35}, minor grid style={gray!18},
    minor tick num=1,
    tick label style={font=\scriptsize, fill=white, fill opacity=0.85, text opacity=1, inner sep=1pt},
    xlabel=$x$, ylabel=$y$,
    xlabel style={anchor=west, xshift=0.25cm}, ylabel style={anchor=south, yshift=0.25cm},
    axis line style={-latex, shorten >=-0.25cm},
    clip mode=individual,
  },
  % A graph that models a situation: axes from zero at the bottom left,
  % labelled with what they measure and its units, no arrowheads on the line.
  contextgraph/.style={
    set layers=labels on top,
    axis lines=left, grid=both,
    major grid style={gray!35}, minor grid style={gray!18},
    tick label style={font=\scriptsize},
    label style={font=\small},
    axis line style={-latex, shorten >=-0.25cm},
    clip mode=individual,
  },
  exampleline/.style={blue!70!black, thick, samples=2, latex-latex},
  lineA/.style={blue!70!black, thick, samples=2},
  lineB/.style={red!75!black, thick, samples=2},
  lineC/.style={green!50!black, thick, samples=2},
  lineD/.style={orange!85!black, thick, samples=2},
  % a legend whose keys show a plain line, without the arrowheads the lines carry
  % It sits to the right of the axes and clear of the x label (0.9 cm past
  % the grid's right edge), with its top level with the grid's top.
  answerlegend/.style={legend cell align=left,
    legend style={at={(1,1)}, anchor=north west, xshift=0.9cm, font=\small, draw=gray!50},
    legend image post style={arrows=-}},
}
% The run and the rise of a slope, drawn on a graph with \draw: dashed, with
% a word beside each.
\tikzset{slopestep/.style={teal, thick, dashed, -latex}}

% ---------- End of the starter's preamble ----------
\renewcommand{\handouttitle}{Pattern: Worked Examples}

\begin{document}

{\LARGE\bfseries\sffamily \doctitle}\par
\medskip
{\color{why} Seven ways to set out a worked example. Copy the one that fits.}

% ---------- A. Substitute and evaluate ----------
\ifanswers
\strand{Substitute into a formula}
\else
\begin{example}{Substitute into a Formula}
The area of a trapezoid is $A = \dfrac{(a + b)h}{2}$. Find the area when
$a = 5$\,cm, $b = 9$\,cm and $h = 4$\,cm.
\begin{align*}
  A &= \frac{(a + b)h}{2} \\
    &= \frac{(\sv{5} + \sv{9})(\sv{4})}{2} \\
    &= \frac{(14)(4)}{2} \\
    &= 28
\end{align*}
The area is 28\,cm$^2$.
\end{example}
\fi

% ---------- B. A method with side notes ----------
\ifanswers
\strand{Expand and simplify}
\else
\begin{example}{Expand and Simplify}
Simplify $2(x + 3) - (x - 4)$.
\begin{alignat*}{2}
  & 2(x + 3) - (x - 4) \qquad && \sidenote{Use the distributive property to expand.} \\
  ={} & 2(x) + 2(3) - 1(x) - 1(-4) \qquad && \sidenote{The minus in front multiplies both terms.} \\
  ={} & 2x + 6 - x + 4 \qquad && \sidenote{Collect like terms.} \\
  ={} & x + 10 &&
\end{alignat*}
\method{Check with a number}
When $x = \sv{1}$: $2(\sv{1} + 3) - (\sv{1} - 4) = 8 + 3 = 11$, and $\sv{1} + 10 = 11$.
\end{example}
\fi

% ---------- C. A table and a graph inside the box ----------
% y = x^2 - 2x - 3 on a grid from x = -3 to 5 and y = -5 to 6: it meets the
% top edge y = 6 at x = 1 - sqrt(10) and 1 + sqrt(10), about -2.16 and 4.16,
% both inside -3 to 5, so the domain stops just short of those.
\ifanswers
\strand{Graph a parabola from a table of values}
\else
\begin{example}{Graph a Parabola from a Table of Values}
\begin{minipage}[t]{0.5\linewidth}
Graph $y = x^2 - 2x - 3$.

\smallskip
Find $y$ for each $x$. When $x = \sv{-1}$,
\[ y = (\sv{-1})^2 - 2(\sv{-1}) - 3 = 0. \]

\medskip
\begin{center}
{\setvaluecol{$-4$}%
\begin{NiceTabular}{N{$2$}N{$5$}}[hvlines]
\CodeBefore \rowcolor{tablehead}{1} \Body
\multicolumn{1}{c}{$x$} & \multicolumn{1}{c}{$y$} \\
$-2$ & $5$ \\
$-1$ & $0$ \\
$0$ & $-3$ \\
$1$ & $-4$ \\
$2$ & $-3$ \\
$3$ & $0$ \\
$4$ & $5$ \\
\end{NiceTabular}}
\end{center}

\medskip
Plot the points. Join them with a smooth curve, to the edge of the grid.
\end{minipage}\hfill
\begin{minipage}[t]{0.46\linewidth}
\vspace{0pt}\centering
\begin{tikzpicture}\begin{axis}[gridgraph, x=0.42cm, y=0.42cm,
    xmin=-3, xmax=5, ymin=-5, ymax=6, xtick={-2,2,4}, ytick={-4,-2,2,4}]
  \addplot[blue!70!black, thick, <->, >=latex, domain=-2.15:4.15, samples=60]{x^2 - 2*x - 3};
  \addplot[only marks, mark size=1.6pt, blue!70!black] coordinates {(-2,5) (-1,0) (0,-3) (1,-4) (2,-3) (3,0) (4,5)};
  \node[anchor=west, font=\small, fill=peach, inner sep=1pt] at (axis cs:4.5,2.5) {$y = x^2 - 2x - 3$};
\end{axis}\end{tikzpicture}
\end{minipage}
\end{example}
\fi

% ---------- D. Two workings side by side ----------
\ifanswers
\strand{Find both intercepts}
\else
\begin{example}{Find Both Intercepts}
Find the intercepts of $2x + 5y = 10$.

\smallskip
\begin{minipage}[t]{0.47\linewidth}\raggedright
At the $x$-intercept, $y = 0$.
\begin{align*}
  2x + 5(\sv{0}) &= 10 \\
  2x &= 10 \\
  x &= 5
\end{align*}
The $x$-intercept is $5$: the point $(5, 0)$.
\end{minipage}\hfill
\begin{minipage}[t]{0.47\linewidth}\raggedright
At the $y$-intercept, $x = 0$.
\begin{align*}
  2(\sv{0}) + 5y &= 10 \\
  5y &= 10 \\
  y &= 2
\end{align*}
The $y$-intercept is $2$: the point $(0, 2)$.
\end{minipage}
\end{example}
\fi

% ---------- G. A slope drawn on the graph: the run and the rise ----------
% y = (3/4)x - 1 on x from -3 to 7 and y from -4 to 5: it meets the left
% edge at y = -3.25 and the right edge at y = 4.25, before the bottom and
% top. From the y-intercept (0, -1): run 4 to the right, then rise 3.
\ifanswers
\strand{Graph a line from its slope and y-intercept}
\else
\begin{example}{Graph a Line from its Slope and $y$-Intercept}
\begin{minipage}[t]{0.52\linewidth}
Graph the line $y = \dfrac34x - 1$.

\smallskip
The slope is $\dfrac34$, so $\dfrac{\text{rise}}{\text{run}} = \dfrac{3}{4}$.

\smallskip
The $y$-intercept is $-1$, so the point $(0, \sv{-1})$ is on the line.

\smallskip
Start at $(0, -1)$. Run $4$ to the right, then rise $3$, to reach
$(4, 2)$. Draw a line through the two points, to the edge of the grid.
\end{minipage}\hfill
\begin{minipage}[t]{0.44\linewidth}
\vspace{0pt}\centering
\begin{tikzpicture}\begin{axis}[gridgraph, x=0.42cm, y=0.42cm,
    xmin=-3, xmax=7, ymin=-4, ymax=5, xtick={-2,2,4,6}, ytick={-2,2,4}]
  \addplot[exampleline, domain=-2.9:6.9]{3/4*x - 1};
  \draw[slopestep] (axis cs:0,-1) -- node[below, font=\scriptsize\sffamily]{run $4$} (axis cs:4,-1);
  \draw[slopestep] (axis cs:4,-1) -- node[right, font=\scriptsize\sffamily]{rise $3$} (axis cs:4,2);
  \addplot[only marks, mark size=1.6pt, blue!70!black] coordinates {(0,-1) (4,2)};
  \node[anchor=west, font=\small, fill=peach, inner sep=1pt] at (axis cs:-2.8,3.6) {$y = \tfrac34x - 1$};
\end{axis}\end{tikzpicture}
\end{minipage}
\end{example}
\fi

% ---------- E. Instructions before a group, and the key ----------
\onlyq{%
\begin{groupnote}
\textbf{\sffamily For questions 1 to 3:} use the method in ``Find Both Intercepts''. Show both workings.
\end{groupnote}}

\begin{questions}
\q{1} \onlyq{Find the $x$-intercept and the $y$-intercept of each line.}
\begin{tasks}(3)
  \task \qa{$3x + 2y = 12$}{$(4, 0)$ and $(0, 6)$}
  \task \qa{$x - 4y = 8$}{$(8, 0)$ and $(0, -2)$}
  \task \qa{$5x + y = -10$}{$(-2, 0)$ and $(0, -10)$}
\end{tasks}
\end{questions}


% ---------- F. A framed note beside an example ----------
\ifanswers
\strand{Percent}
\else
\begin{example}{Percent of an Amount}
\begin{minipage}[t]{0.66\linewidth}
A jacket costs \$60. It is on sale for 15\% off. How much is taken off?
\begin{align*}
  15\% \text{ of } \$60 &= 0.15 \times \$60 \\
  &= \$9
\end{align*}
The sale takes \$9 off the price.
\end{minipage}\hfill
\begin{minipage}[t]{0.3\linewidth}
\vspace{0pt}
\begin{notebox}
$15\%$ means $\dfrac{15}{100}$, or $0.15$.
\end{notebox}
\end{minipage}
\end{example}
\fi

\end{document}
